\BourbakiExercisesSection

\begin{enumerate}
\def\labelenumi{\arabic{enumi}.}
\item
  If V is a finitely generated projective A-module, so is End(V), canonically identified with \(V \otimes_{A} V^{*}\) (II, §4, no. 2, Corollary to Proposition 2). If every element \(u \in \text{End}(V)\) is associated with the A-linear form \(v \mapsto \text{Tr}(uv)\) on End(V), an A-linear bijection is defined of End(V) onto its dual (End(V))*. Identifying End(V) and its dual A-module under this mapping, the A-algebra structure on End(V) defines by duality (no. 1, Example 3) an A-cogebra structure on End(V), which is coassociative and has as counit the form Tr: \(V \to A\) . In particular, for \(V = A^{n}\) , a cogebra structure is defined on \(\mathbf{M}_{n}(\mathbf{A})\) for which the coproduct is given by \(c(E_{ij}) = \sum_{k} E_{kj} \otimes E_{ik}\) . This cogebra structure and the algebra structure on \(\mathbf{M}_{n}(\mathbf{A})\) do not define a bigebra structure.
\item
  Show that in Definition 3 (no. 4), condition (3) (relative to graded bigebras) is equivalent to the following: the product \(m: \mathbf{E} \otimes \mathbf{E} \to \mathbf{E}\) is a morphism of the graded cogebra \(\mathbf{E} \otimes \mathbf{E}\) into the graded cogebra \(\mathbf{E}\) .
\item
  Show that if \(\mathbf{M}\) is an A-module there exists on \(\mathsf{T}(\mathbf{M})\) one and only one cocommutative bigebra structure whose algebra structure is the usual structure and whose coproduct \(c\) is such that, for all \(x\in \mathbf{M}\) , \(c(x) = 1\otimes x + x\otimes 1\) .
\item
  Let \(\mathbf{E}\) be a commutative bigebra over \(\mathbf{A}\) , \(m\) the product \(\mathbf{E} \otimes \mathbf{E} \to \mathbf{E}\) , \(c\) the coproduct \(\mathbf{E} \to \mathbf{E} \otimes \mathbf{E}\) , \(e\) the unit element and \(\gamma\) the counit of \(\mathbf{E}\) .
\end{enumerate}

\begin{enumerate}
\def\labelenumi{(\alph{enumi})}
\tightlist
\item
  Show that for every commutative A-algebra B, the law of composition which associates with every ordered pair
\end{enumerate}

\[
(u, v) \in \operatorname{Hom} _ {\mathbf {A} - \mathrm{alg.}} (\mathbf {E}, \mathbf {B}) \times \operatorname{Hom} _ {\mathbf {A} - \mathrm{alg.}} (\mathbf {E}, \mathbf {B})
\]

the A-algebra homomorphism

\[
\mathrm{E} \xrightarrow {c} \mathrm{E} \otimes \mathrm{E} \xrightarrow {u \otimes v} \mathrm{B} \otimes \mathrm{B} \xrightarrow {m _ {\mathrm{B}}} \mathrm{B}
\]

(where \(m_{\mathbf{B}}\) is the product in B) defines a monoid structure on \(\operatorname{Hom}_{\mathbf{A}\text{-alg.}}(\mathbf{E}, \mathbf{B})\) . (b) In order that, for every commutative A-algebra B, \(\operatorname{Hom}_{\mathbf{A}\text{-alg.}}(\mathbf{E}, \mathbf{B})\) be a group, it is necessary and sufficient that there exist an A-algebra homomorphism \(i: \mathbf{E} \to \mathbf{E}\) such that \(i(e) = e\) and the composite mappings \(m \circ (1_{\mathbf{E}} \otimes i) \circ c\) and \(m \circ (i \otimes 1_{\mathbf{E}}) \circ c\) are both equal to the mapping \(x \mapsto \gamma(x)e\) . \(i\) is then called an inversion in E; it is an isomorphism of E onto the opposite bigebra \(E^0\) .

\begin{enumerate}
\def\labelenumi{\arabic{enumi}.}
\setcounter{enumi}{4}
\item
  Let \(E_{1}\) , \(E_{2}\) be two bigebras over A. Show that on the tensor product \(E_{1} \otimes_{A} E_{2}\) the algebra and cogebra structures the tensor product of those on \(E_{1}\) and \(E_{2}\) define on \(E_{1} \otimes_{A} E_{2}\) a bigebra structure.
\item
  Let M be a graded A-module of type N. Define a skew graded bigebra structure on the universal anticommutative algebra G(M) (§ 7, Exercise
\end{enumerate}

13); deduce a canonical graded algebra homomorphism \(\mathbf{G}(\mathbf{M}^{*}) \to \mathbf{G}(\mathbf{M})^{*\mathbf{gr}}\) and notions of inner product between elements of \(\mathbf{G}(\mathbf{M})\) and elements of \(\mathbf{G}(\mathbf{M}^{*})\) .

\begin{enumerate}
\def\labelenumi{\arabic{enumi}.}
\setcounter{enumi}{6}
\tightlist
\item
  With the hypotheses and notation of Proposition 11 (no. 9), prove the formula
\end{enumerate}

\[
(\det g) \phi^ {- 1} = \bigwedge (g) \circ \phi^ {- 1} \circ \bigwedge (^ {t} g).
\]

Form this formula and (79) (no. 11) derive a new proof of the expression for the inverse of a square matrix using the matrix of its cofactors (§ 8, no. 6, formula (26)).

\begin{enumerate}
\def\labelenumi{\arabic{enumi}.}
\setcounter{enumi}{7}
\tightlist
\item
  Let E be a free A-module of dimension n.~For every A-linear mapping v of \(\bigwedge(\mathbf{M})\) into an A-module G, we can write \(v = \sum_{p=0}^{n} v_{p}\) , where \(v_{p}\) is the restriction of v to \(\bigwedge^{p}(\mathbf{M})\) ; we then set \(\eta v = \sum_{p=0}^{n} (-1)^{p} v_{p}\) , so that \(\eta^{2} v = v\) . Let u be an isomorphism of M onto its dual \(M^{*}\) and let \(\Delta\) be the determinant of the matrix of u with respect to a basis \((e_{i})\) of M and the dual basis of \(M^{*}\) (a determinant which does not depend on the basis of M chosen). If \(\phi\) is the isomorphism of \(\bigwedge(\mathbf{M})\) onto \(\bigwedge(\mathbf{M}^{*})\) defined starting with \(e = e_{1} \wedge e_{2} \wedge \cdots \wedge e_{n}\) , show the formula
\end{enumerate}

\[
\eta^ {n + 1} \Delta \phi = \bigwedge (t u) \circ \phi^ {- 1} \circ \bigwedge (u).
\]

\begin{enumerate}
\def\labelenumi{\arabic{enumi}.}
\setcounter{enumi}{8}
\tightlist
\item
  The adjoint of a square matrix \(X\) of order \(n\) over \(A\) is the matrix \(\tilde{X} = (\det(X^{it}))\) of minors of \(X\) of order \(n - 1\) . Show that if \(X\) is invertible then \(\det(\tilde{X}) = (\det X)^{n-1}\) and that every minor \(\det(\tilde{X}_{H,K})\) of \(\tilde{X}\) of order \(p\) is given by the formula
\end{enumerate}

\[
\det (\tilde {X} _ {\mathrm{H}, \mathrm{K}}) = (\det X) ^ {p - 1} (X _ {\mathrm{H} ^ {\prime}, \mathrm{K} ^ {\prime}})\tag{1}
\]

where H' and K' are the complements of H and K respectively in \([1, n]\) (``Jacobi's identities''; use relation (80) of no. 11).

\begin{enumerate}
\def\labelenumi{\arabic{enumi}.}
\setcounter{enumi}{9}
\tightlist
\item
  From every identity \(\Phi = 0\) between minors of an arbitrary invertible square matrix \(X\) of order \(n\) over A another identity \(\tilde{\Phi} = 0\) can be derived called the complement of \(\Phi = 0\) , by applying the identity \(\Phi = 0\) to the minors of the adjoint \(\tilde{X}\) of \(X\) and then replacing the minors of \(\tilde{X}\) by functions of the minors of \(X\) using identity (1) of Exercise 9. In this way show the following identity:
\end{enumerate}

\[
\det (X ^ {i h}) \det (X ^ {j k}) - \det (X ^ {i k}) \det (X ^ {j h}) = (\det X) \det (X ^ {i j, h k}) \quad \text { for } \quad i <   j
\]

where \(X^{ij, hk}\) denotes the minor of X of order n - 2 obtained by suppressing in X the rows of index i, j and the columns of index h, k.

¶ 11. Let \(\Phi = 0\) be an identity between minors of an invertible square matrix \(X\) of order \(n\) over a commutative ring A, \(\tilde{\Phi} = 0\) the complementary identity (Exercise 10) and \(k\) an integer \(>0\) . Let \(Y\) be an invertible square matrix of order \(n + k\) and let \(\tilde{Y}_0\) be the submatrix of the adjoint matrix \(\tilde{Y}\) formed by suppressing in \(\tilde{Y}\) the rows of index \(\leqslant k\) and the columns of index \(\leqslant k\) . If \(\tilde{Y}_0\) is assumed to be invertible and the identity \(\tilde{\Phi} = 0\) is applied to the minors of \(\tilde{Y}_0\) and then each minor which appears in this identity (considered as a minor of \(\tilde{Y}\) ) is replaced by its expression as a function of the minors of \(Y\) , using identity (1) of Exercise 9, an identity \(\Phi_k = 0\) is obtained between minors of the matrix \(Y\) (valid if \(Y\) and \(\tilde{Y}_0\) are invertible) which is called the extension of order \(k\) of the identity \(\Phi = 0\) .

In particular, Let \(A = (\alpha_{ij})\) be an invertible square matrix of order \(n + k\) , B the submatrix of A of order k obtained by suppressing in A the rows and columns of index \textgreater k and \(\Delta_{ij}\) the determinant of the matrix of order \(k + 1\) obtained by suppressing in A the rows of index \textgreater k except that of index \(k + i\) and the columns of index \textgreater k except that of index \(k + j\) ; if C denotes the matrix \((\Delta_{ij})\) of order n, show the identity

\[
\det C = (\det A) (\det B) ^ {n - 1}
\]

valid when \(A\) and \(B\) are invertible (show that it is the extension of order \(k\) of the total expansion of a determinant of order \(n\) ).

State the identities obtained by extending the Laplace expansion (§ 8, no. 6) and the identity of Exercise 2 of § 8.

¶ 12. Let \(X = (\xi_{ij})\) be a square matrix of order \(n\) over a commutative field \(K\) ; let \(H\) be a subset of \([1, n]\) with \(p\) elements and \(H'\) the complement of \(H\) in \([1, n]\) ; suppose that, for every ordered pair of indices \(h \in H\) , \(k \in H'\) ,

\[
\sum_ {i} \xi_ {i h} \xi_ {i k} = 0.
\]

Show that, for all subsets L, M of \([1, n]\) with p elements,

\[
\rho_ {\mathbf {M}, \mathbf {M} ^ {\prime}} \det (X _ {\mathrm{L}, \mathrm{H}}) \det (X _ {\mathbf {M} ^ {\prime}, \mathbf {H} ^ {\prime}}) - \rho_ {\mathrm{L}, \mathrm{L} ^ {\prime}} \det (X _ {\mathbf {M}, \mathbf {H}}) \det (X _ {\mathrm{L}, \mathrm{H} ^ {\prime}}) = 0.
\]

(Consider the columns \(x_h\) of index \(h \in \mathbf{H}\) in \(X\) as vectors in \(\mathbf{E} = \mathbf{K}^n\) and the columns \(x_k^*\) of index \(k \in \mathbf{H}'\) as vectors in \(\mathbf{E}^*\) and show that the \((n - p)\) -form \(x_{\mathbf{H}'}^*\) is proportional to \(\phi_p(x_{\mathbf{H}})\) .)

\begin{enumerate}
\def\labelenumi{\arabic{enumi}.}
\setcounter{enumi}{12}
\tightlist
\item
  Let \(\Gamma\) and \(\Delta\) be two determinants of invertible matrices of order n over a commutative ring and \(\Gamma_{ij}\) the determinant obtained by replacing in \(\Gamma\) the column of index i by the column of \(\Delta\) of index j. Show that
\end{enumerate}

\[
\det (\Gamma_ {i j}) = \Gamma^ {n - 1} \Delta .
\]

(Expand \(\Gamma_{ij}\) by the \(i\) -th column and use Exercise 9.)

\begin{enumerate}
\def\labelenumi{\arabic{enumi}.}
\setcounter{enumi}{13}
\item
  Let M be a free A-module of dimension n \textgreater{} 1. Show that every isomorphism \(\psi\) of \(\bigwedge^{p}(M)\) onto \(\bigwedge^{n-p}(M^{*})\) such that \(\bigwedge^{n-p}(\check{u}) \circ \psi = \psi \circ \bigwedge^{p}(u)\) for every automorphism u of M, of determinant equal to 1, is one of the isomorphisms \(\phi_{p}\) defined in Proposition 12 of no. 11 (proceed as in §7, Exercise 4).
\item
  Let E be a free A-module of dimension n.~Let z be a p-vector over E and \(z^{*}\) a \((p + q)\) -form on E; z can be canonically identified (§ 7, Exercise 8) with the skewsymmetrization \(az_{1}\) of a contravariant tensor \(z_{1}\) of order p and \(z^{*}\) with the skewsymmetrization of a covariant tensor of order \(p + q\) . If, in the mixed tensor \(z_{1}\) . \(z^{*}\) , the k-th contravariant index and the \((p + k)\) -th covariant index are contracted for \(1 \leqslant k \leqslant p\) , show that the covariant tensor of order q thus obtained is skewsymmetrized and is canonically identified with the q-form \(z^{*} \perp z\) , up to sign.
\item
  Let E be a free A-module of dimension n.~The regressive product of \(x \in \bigwedge(\mathrm{E})\) and \(y \in \bigwedge(\mathrm{E})\) with respect to a basis \(\{e\}\) of \(\bigwedge^{n}(\mathrm{E})\) , denoted by \(x \vee y\) , is the element \(\phi(\phi(x) \wedge \phi(y))\) , the isomorphism \(\phi\) being relative to e. This product is only defined to within an invertible factor, depending on the basis \(\{e\}\) chosen. Show that if x is a p-vector and y a q-vector, \(x \vee y = 0\) if \(p + q < n\) and that, if \(p + q \geqslant n\) , \(x \vee y\) is a \((p + q - n)\) -vector such that
\end{enumerate}

\[
y \vee x = (- 1) ^ {(n - p) (n - q)} x \vee y.
\]

The regressive product is associative and distributive with respect to addition and defines on \(\bigwedge(\mathrm{E})\) a unital algebra structure isomorphic to the exterior algebra structure. Express the components of \(x \vee y\) as a function of those of \(x\) and \(y\) for a given basis on \(\mathbf{E}\) .

\begin{enumerate}
\def\labelenumi{\arabic{enumi}.}
\setcounter{enumi}{16}
\tightlist
\item
  Let K be a commutative field and E a vector space over K.
\end{enumerate}

\begin{enumerate}
\def\labelenumi{(\alph{enumi})}
\item
  Let \(z\) be a \(p\) -vector \(\neq 0\) over \(E\) and let \(V_{z}\) be the vector subspace of \(E\) consisting of the vectors \(x\) such that \(z \wedge x = 0\) . Then \(\dim(V_{z}) \leqslant p\) ; if \((x_{i})_{1 \leqslant i \leqslant q}\) is a free system of vectors of \(V_{z}\) , there exists a \((p - q)\) -vector \(v\) such that \(z = v \wedge x_{1} \wedge \dots \wedge x_{q}\) . If \(E\) is finite-dimensional, then \(V_{z} = \phi(M_{z}^{0})\) , where \(M_{z}^{0}\) is the orthogonal in \(E^{*}\) of the subspace \(M_{z}\) associated with \(z\) . For \(z\) to be a pure \(p\) -vector, it is necessary and sufficient that \(\dim(V_{z}) = p\) and then \(V_{z} = M_{z}\) .
\item
  Let \(v\) be a pure \(p\) -vector, \(w\) a pure \(q\) -vector and \(V\) and \(W\) the subspaces of \(E\) associated respectively with \(v\) and \(w\) . In order that \(V \subset W\) , it is necessary and sufficient that \(q \geqslant p\) and that there exist a \((q - p)\) -vector \(u\) such that \(w = u \wedge v\) . In order that \(V \cap W = \{0\}\) , it is necessary and sufficient that \(v \wedge w \neq 0\) ; the subspace \(V + W\) is then associated with the pure \((p + q)\) -vector \(v \wedge w\) .
\item
  Let \(u\) and \(v\) be two pure \(p\) -vectors and \(U\) and \(V\) the subspaces associated respectively with \(u\) and \(v\) . For \(u + v\) to be pure, it is necessary and sufficient that \(\dim(U \cap V) \geqslant p - 1\) .
\end{enumerate}

\begin{enumerate}
\def\labelenumi{\arabic{enumi}.}
\setcounter{enumi}{17}
\item
  Let \(z = \sum_{H} \alpha_{H} e_{H}\) be a non-zero pure p-vector of an n-dimensional vector space E, expressed using its components with respect to an arbitrary basis \((e_{i})_{1 \leqslant i \leqslant n}\) of E. Let G be a subset of \([1, n]\) with p elements such that \(\alpha_{G} \neq 0\) ; let \((i_{h})_{1 \leqslant h \leqslant p}\) be the sequence of indices in G arranged in increasing order and \((j_{k})_{1 \leqslant k \leqslant n - p}\) the sequence of indices in the complement \(G'\) of G, arranged in increasing order. For every ordered pair \((h, k)\) of indices such that \(1 \leqslant h \leqslant p\) , \(1 \leqslant k \leqslant n - p\) , let \(\beta_{hk}\) be the component \(\alpha_{H}\) of z corresponding to the subset H of \([1, n]\) consisting of the p - 1 indices in G distinct from \(i_{h}\) and of index \(j_{k} \in G'\) ; let X be the matrix \((\beta_{hk})\) with p rows and n - p columns. Given an arbitrary subset L of \([1, n]\) with p elements such that \(L \cap G\) has \(q \geqslant 1\) elements, show that \((\alpha_{G})^{q-1} \alpha_{L}\) is equal, up to sign, to the minor of X of order q, consisting of the rows of index h such that \(i_{h} \in G \cap G\) and the columns of index k such that \(j_{k} \in L \cap G\) . (Write z as being of the form \(\alpha_{G} y_{1} \wedge y_{2} \wedge \cdots \wedge y_{p}\) , where the vectors \(y_{i}\) are such that, in the matrix Y with n rows and p columns whose columns are the \(y_{i}\) , the submatrix consisting of the rows whose indices belong to G is the unit matrix.)
\item
  Over a finite-dimensional vector space E, let \(z = x_{1} \wedge x_{2} \wedge \cdots \wedge x_{p}\) be a pure p-vector \(\neq 0\) ; let \(v^{*}\) be a q-form ( \(q \leqslant p\) ) and \(u^{*}\) an arbitrary element of \(\bigwedge(\mathrm{E}^{*})\) . Show that
\end{enumerate}

\[
(u ^ {*} \wedge v ^ {*}) \lrcorner z = (- 1) ^ {q (q - 1) / 2} \sum_ {\mathrm{H}} \rho_ {\mathrm{H}, \mathrm{K}} \langle v ^ {*}, x _ {\mathrm{H}} \rangle (u ^ {*} \lrcorner x _ {\mathrm{K}})
\]

where H runs through the set of subsets of \([1, p]\) with q elements, K is the complement of H in \([1, p]\) and \(x_{H}\) (resp. \(x_{K}\) ) denotes the exterior product of the \(x_{i}\) such that \(i \in H\) (resp. \(i \in K\) ).

\begin{enumerate}
\def\labelenumi{\arabic{enumi}.}
\setcounter{enumi}{19}
\item
  Let E be an n-dimensional vector space, z a pure p-vector over E, V the vector subspace of E associated with z, \(u^{*}\) a pure q-vector over \(E^{*}\) , \(W'\) the subspace of \(E^{*}\) associated with \(u^{*}\) and W the subspace of E, of dimension n - q, orthogonal to \(W'\) . If q \textless{} p, \(u^{*} \lrcorner z\) is a pure (p - q)-vector over E, which is zero if \(\dim(V \cap W) > p - q\) ; if \(\dim(V \cap W) = p - q\) , \(V \cap W\) is associated with \(u^{*} \lrcorner z\) .
\item
  \begin{enumerate}
  \def\labelenumii{(\alph{enumii})}
  \tightlist
  \item
    Show directly (without using the isomorphisms \(\phi_p\) ) that every \((n - 1)\) -vector over an \(n\) -dimensional vector space is pure.
  \end{enumerate}
\end{enumerate}

\begin{enumerate}
\def\labelenumi{(\alph{enumi})}
\setcounter{enumi}{1}
\tightlist
\item
  Let A be a commutative algebra over a field K with a basis consisting of the unit element 1 and three elements \(a_{1}\) , \(a_{2}\) , \(a_{3}\) whose products with one another are zero. Let E be the A-module \(\mathbf{A}^3\) and \((e_i)_{1\leqslant i\leqslant 3}\) its canonical basis; show that the bivector
\end{enumerate}

\[
a _ {1} (e _ {2} \land e _ {3}) + a _ {2} (e _ {3} \land e _ {1}) + a _ {3} (e _ {1} \land e _ {2})
\]

is not pure.

\begin{enumerate}
\def\labelenumi{\arabic{enumi}.}
\setcounter{enumi}{21}
\tightlist
\item
  Let E be an ordered set in which every interval \([a, b]\) is finite. Let S be the subset of \(E \times E\) consisting of the ordered pairs \((x, y)\) such that \(x \leqslant y\) . Show that a coassociative cogebra structure is defined on the A-module \(\mathbf{A}^{(\mathrm{s})}\) of formal linear combinations of the elements of S by taking as coproduct
\end{enumerate}

\[
c ((x, y)) = \sum_ {x \leqslant z \leqslant y} (x, z) \otimes (z, y).
\]

For this cogebra the linear form \(\gamma\) such that

\[
\gamma ((x, y)) = 0 \quad \text { if } \quad x \neq y \text { in   E }, \quad \gamma ((x, x)) = 1
\]

is a counit.

\begin{enumerate}
\def\labelenumi{\arabic{enumi}.}
\setcounter{enumi}{22}
\tightlist
\item
  Let C be a cogebra over a ring A. A sub-A-module S of C is called a subcogebra of C if \(c(S) \subset \operatorname{Im}(S \otimes S)\) .
\end{enumerate}

\begin{enumerate}
\def\labelenumi{(\alph{enumi})}
\item
  If C is coassociative (resp. cocommutative), so is every subcogebra. If \(\gamma\) is a counit of C, the restriction of \(\gamma\) to any subcogebra of C is a counit.
\item
  Which are the subcogebras of the cogebra \(\mathbf{A}^{(\mathbf{X})}\) (no. 1, Example 4).?
\item
  If \((\mathbf{S}_{\alpha})\) is a family of subcogebras of \(\mathbf{C}\) , \(\sum_{\alpha} \mathbf{S}_{\alpha}\) is a subcogebra of \(\mathbf{C}\) (the smallest containing all the \(\mathbf{S}_{\alpha}\) ).
\end{enumerate}

\begin{enumerate}
\def\labelenumi{\arabic{enumi}.}
\setcounter{enumi}{23}
\tightlist
\item
  Let C be a cogebra over a ring A. A sub-A-module S of C is called a right (resp. left) coideal if \(c(S) \subset \operatorname{Im}(S \otimes C)\) (resp. \(c(S) \subset \operatorname{Im}(C \otimes S)\) ).
\end{enumerate}

\begin{enumerate}
\def\labelenumi{(\alph{enumi})}
\item
  A subcogebra (Exercise 23) is a right and left coideal. The converse is true if \(\mathbf{A}\) is a field.
\item
  Every sum of right (resp. left) coideals is a right (resp. left) coideal.
\item
  If S is a right (resp. left) coideal of C, the orthogonal \(S^{0}\) is a right (resp. left) ideal of the dual algebra \(C^{*}\) .
\item
  If A is a field and \(\mathfrak{I}'\) a right (resp. left) ideal of the dual algebra \(\mathbf{C}^*\) , the orthogonal \(\mathfrak{I}''^0\) in \(\mathbf{C}\) is a right (resp. left) coideal of \(\mathbf{C}\) . (Argue by reductio ad absurdum by assuming that there exists an \(x \in \mathfrak{I}''^0\) such that \(c(x) \notin \mathfrak{I}''^0 \otimes \mathbf{C}\) ; write \(c(x) = \sum_{i} y_i \otimes z_i\) , where the \(z_i\) are linearly independent, and show that there would be \(y' \in \mathfrak{I}'\) and \(z' \in \mathbf{C}^*\) such that \(\langle y' z', x \rangle \neq 0\) .)
\end{enumerate}

Deduce that for a vector subspace S of C to be a right (resp. left) coideal, it is necessary and sufficient that \(S^{0}\) be a right (resp. left) ideal of \(C^{*}\) .

\begin{enumerate}
\def\labelenumi{(\alph{enumi})}
\setcounter{enumi}{4}
\tightlist
\item
  Deduce from (d) that if A is a field every intersection of right coideals (resp. left coideals, resp. subcogebras) of C is a right coideal (resp. left coideal, resp. subcogebra). For S to be a subcogebra of C, it is necessary and sufficient that its orthogonal \(S^{0}\) be a two-sided ideal of the algebra \(C^{*}\) .
\end{enumerate}

\begin{enumerate}
\def\labelenumi{\arabic{enumi}.}
\setcounter{enumi}{24}
\tightlist
\item
  Let C be a cogebra over a ring A. A sub-A-module S of C is called a coideal if \(c(\mathbf{S}) \subset \operatorname{Im}(\mathbf{S} \otimes \mathbf{C}) + \operatorname{Im}(\mathbf{C} \otimes \mathbf{S})\) . If C is counital, with counit \(\gamma\) , a coideal S is called conull if \(\gamma(x) = 0\) for \(x \in S\) .
\end{enumerate}

\begin{enumerate}
\def\labelenumi{(\alph{enumi})}
\item
  A right or left coideal is a coideal. Every sum of coideals is a coideal.
\item
  If S is a coideal of C, the orthogonal \(S^{0}\) is a subalgebra of the dual algebra \(C^{*}\) . If C is counital and S is conull, \(S^{0}\) is unital.
\item
  Suppose that A is a field and C a counital cogebra. Show that if \(E'\) is a unital subalgebra of \(C^{*}\) with the same unit element as \(C^{*}\) , the orthogonal \(E'^{0}\) in C is a conull coideal.
\item
  If S is a coideal of C, show that there exists on C/S one and only one cogebra structure such that the canonical mapping \(p: C \rightarrow C/S\) is a cogebra morphism. If C is counital and S is conull, C/S is counital.
\item
  For every cogebra morphism \(f\colon\mathbf{C}\to\mathbf{C}^{\prime},\mathbf{S}=\operatorname{Ker}(f)\) is a coideal of \(\mathbf{C}\) , \(\mathbf{S}^{\prime}=\operatorname{Im}(f)\) a subcogebra of \(\mathbf{C}^{\prime}\) and, in the canonical factorization
\end{enumerate}

\[
\mathrm{C} \xrightarrow {p} \mathrm{C} / \mathrm{S} \xrightarrow {g} \mathrm{S} ^ {\prime} \xrightarrow {j} \mathrm{C} ^ {\prime}
\]

of \(f, g\) is a cogebra isomorphism.

\begin{enumerate}
\def\labelenumi{(\arabic{enumi})}
\setcounter{enumi}{25}
\tightlist
\item
  Let C be a coassociative counital cogebra over a commutative field K.
\end{enumerate}

\begin{enumerate}
\def\labelenumi{(\alph{enumi})}
\item
  For every vector subspace E of C which is a left C*-module under the law \((u, x) \mapsto u \perp x\) , the left annihilator of this C*-module is a two-sided ideal \(S'\) of the algebra \(C^*\) . Show that if E is finite-dimensional over K, \(C^*/S'\) is finite-dimensional.
\item
  Deduce from (a) that for every element \(x \in \mathbb{C}\) , the subcogebra of \(\mathbb{C}\) generated by \(x\) is finite-dimensional. (Note that the \(\mathbb{C}^*\) -module \(\mathbb{E}\) generated by \(x\) is finite-dimensional and that, if \(\mathfrak{S}'\) is its left annihilator, \(x \in \mathfrak{S}'^0\) , the orthogonal of \(\mathfrak{S}'\) in \(\mathbb{C}\) .)
\end{enumerate}

\begin{enumerate}
\def\labelenumi{(\arabic{enumi})}
\setcounter{enumi}{26}
\tightlist
\item
  Let B be an associative unital algebra over a commutative field K and let \(m: B \otimes_{K} B \to B\) be the K-linear mapping defining the multiplication on B. The tensor product \(B^{*} \otimes_{K} B^{*}\) (where \(B^{*}\) is the dual vector space of the vector space B) is canonically identified with a vector subspace of \((\mathbf{B} \otimes_{\mathbf{K}} \mathbf{B})^{*}\) (II, §7, no. 7, Proposition 16).
\end{enumerate}

\begin{enumerate}
\def\labelenumi{(\alph{enumi})}
\item
  Show that for an element \(w \in \mathbf{B}^*\) the following conditions are equivalent: (1) \({}^t m(w) \in \mathbf{B}^* \otimes \mathbf{B}^*\) ; (2) the set of \(x \lrcorner w\) , where \(x\) runs through \(\mathbf{B}\) , is a finite-dimensional subspace of \(\mathbf{B}^*\) ; (3) the set of \(w \lfloor x\) , where \(x\) runs through \(\mathbf{B}\) , is a finite-dimensional subspace of \(\mathbf{B}^*\) ; (4) the set of \(x \lrcorner w \lfloor y\) , where \(x\) and \(y\) run through \(\mathbf{B}\) , is contained in a finite-dimensional subspace of \(\mathbf{B}^*\) .
\item
  Let \(\mathbf{B}' \subset \mathbf{B}^*\) be the set of \(w \in \mathbf{B}^*\) satisfying the equivalent conditions in
\end{enumerate}

(a). Show that \({}^t m(\mathbf{B}') \subset \mathbf{B}' \otimes_{\mathbb{K}} \mathbf{B}'\) . (For \(w \in \mathbf{B}'\) , write \({}^t m(w) = \sum_i u_i \otimes v_i\) , where for example the \(u_i\) are linearly independent in \(\mathbf{B}^*\) ; show that then \(v_i \in \mathbf{B}'\) for all \(i\) using the associativity of \(\mathbf{B}\) and arguing by reductio ad absurdum.) \(\mathbf{B}'\) is therefore the largest cogebra contained in \(\mathbf{B}^*\) , for which \({}^t m\) is the coproduct. \(\mathbf{B}'\) is called the dual cogebra of the algebra \(\mathbf{B}\) . When \(\mathbf{B}\) is finite-dimensional, \(\mathbf{B}' = \mathbf{B}^*\) .

\begin{enumerate}
\def\labelenumi{(\alph{enumi})}
\setcounter{enumi}{2}
\item
  Show that if \(\mathbf{B}\) is a commutative field of infinite dimension over \(\mathbf{K}\) , \(\mathbf{B}' = \{0\}\) . (Note that the orthogonal of the set of \(x \perp w\) , for same \(w \in \mathbf{B}^*\) , \(x\) running through \(\mathbf{B}\) , is an ideal of \(\mathbf{B}\) .)
\item
  Let \(\mathbf{B}_1, \mathbf{B}_2\) be two K-algebras and \(f: \mathbf{B}_1 \to \mathbf{B}_2\) a K-homomorphism of algebras. Show that \(^t f(\mathbf{B}_2') \subset \mathbf{B}_1'\) and that \(^t f\) , restricted to \(\mathbf{B}_2'\) , is a cogebra homomorphism of \(\mathbf{B}_2'\) into \(\mathbf{B}_1'\) .
\item
  If \(\mathbf{C}\) is a coassociative counital cogebra over \(\mathbf{K}\) , the canonical injection \(\mathbf{C} \to \mathbf{C}^{**}\) of the vector space \(\mathbf{C}\) into its bidual maps \(\mathbf{C}\) onto a subspace of \((\mathbf{C}^*)'\) , the dual cogebra of the algebra \(\mathbf{C}^*\) dual to \(\mathbf{C}\) , and is a cogebra homomorphism of \(\mathbf{C}\) into \((\mathbf{C}^*)'\) .
\end{enumerate}

